\documentclass[a4paper,11pt,twoside]{article}
\usepackage[utf8]{inputenc}	% Text coding
\usepackage[T1]{fontenc}
\usepackage{lmodern}
\usepackage[czech]{babel}
\usepackage{epsfig}
\usepackage{amsfonts,amsmath,amssymb}
\usepackage{graphicx}
\usepackage[unicode]{hyperref}
\usepackage{indentfirst}
\usepackage{fancyhdr}
\usepackage{xifthen}
\usepackage{amsthm,thmtools}
\usepackage{bold-extra}
\usepackage[dvipsnames]{xcolor}
\usepackage[subrefformat=simple,labelformat=simple]{subcaption} % Instead of subfigure
\usepackage{listings}
\usepackage{comment}
\usepackage{titlesec}
\usepackage{underscore}
\usepackage{makecell}       % Šířky čar v tabulkách
\usepackage{wrapfig}

% Page size
\addtolength{\topmargin}{-1.5cm} %\addtolength{\textheight}{-10cm}
\addtolength{\textwidth}{4cm} \addtolength{\textheight}{4cm} % Width and height of the text
\addtolength{\voffset}{-0.5cm} % Top margin
\addtolength{\hoffset}{-2cm}
\setlength{\headheight}{15pt}

\DeclareMathOperator{\e}{e}

\def\vector#1{\boldsymbol{#1}}								% Vector
\renewcommand{\d}{\mathrm{d}}
\newcommand{\derivative}[3][]{\ifthenelse{\isempty{#1}}	    % Normal derivative
	{\frac{\d{#2}}{\d{#3}}}
	{\frac{\d^{#1}{#2}}{\d{#3}^{#1}}}
}
\newcommand{\im}{\mathrm{i}}

\def\makematrix#1{\begin{pmatrix}#1\end{pmatrix}}       % Matrix
\def\abs#1{\left|#1\right|}
\def\probability#1{\mathrm{Pr}\left[#1\right]}
\def\expectation#1{\mathrm{E}\left[#1\right]}
\def\dispersion#1{\sigma_{#1}^{2}}

\def\code#1{\textnormal{\texttt{#1}}}
\def\file#1{\textnormal{\textbf{\texttt{#1}}}}
\def\ghfile#1#2{\textnormal{\textbf{\texttt{\href{https://github.com/PavelStransky/PCInPhysics2025/blob/main/#1#2}{#2}}}}}

\def\abbreviation#1{\textnormal{\textsc{#1}}}

\begin{document}

\section*{Domácí úkol na 10.3.2025}
\subsection*{Náhodná čísla a histogram}
    Nakreslete \emph{normalizovaný histogram} a vypočítejte výběrovou střední hodnotu a rozptyl pro následující vstupní data:
\begin{enumerate}
    \item
        Výběr z $R(0,1)$, což označuje náhodnou veličinu rovnomětně rozdělenou na intervalu $\langle 0,1\rangle$ (v Pythonu generované například pomocí funkce \code{numpy.random.rand()}).

    \item
        Výběr ze \emph{součtu dvou} rovnoměrných rozdělení na intervalu $\langle 0,1\rangle$.

    \item
        Výběr ze \emph{součtu $M$} rovnoměrných rozdělení na intervalu $\langle 0,1\rangle$.
        Pro $M=12$ porovnejte s hustotou pravděpodobnosti Gausovského normálního rozdělení $N(\mu,\sigma^2)$ danou
        \begin{equation}
            \rho_N(x)=\frac{1}{\sqrt{2\pi\sigma^{2}}}\e^{-\frac{(x-\mu)^{2}}{2\sigma^{2}}},
        \end{equation}
        pro $\mu=6$ a $\sigma=1$.
        Na základě porovnání navrhněte jednoduchý generátor náhodných čísel s Gaussovským rozdělením $N(0,1)$ pomocí rovnoměrného rozdělení $R(0,1)$.

    \item
        Nakreslete histogram hodnot logistického zobrazení
        \begin{equation}
            x_{n+1}=rx_{n}(1-x_{n}),
        \end{equation}	
        pro $r=4$.
        Hodilo by se toto zobrazení ke generování rovnoměrně rozdělených náhodných čísel na intervalu $\langle 0,1\rangle$?

    \item
        Ukažte numericky, že vlastní hodnoty $a_{k}$ velké reálné náhodné symetrické matice
        \begin{equation}
            M=\frac{1}{2}\left(A+A^{\mathrm{T}}\right),
        \end{equation}	
        kde $A_{jk}\in N(0,\sigma^{2})$ a $A^{\mathrm{T}}$ je matice transponovaná k $A$, má rozdělení vlastních hodnot přibližně rovné Wignerovu polokruhovému zákonu
        \begin{equation}
            \rho_{W}(a)=\frac{2}{\pi R^{2}}\left(R^{2}-a^{2}\right),\qquad R=2\sigma\sqrt{N},
        \end{equation}

\end{enumerate}
Pro pěkné grafy volte alespoň $N=10000$ (počet prvků výběru) a $H=100$ (počet intervalů histogramu).

\end{document}
